\subsection{Scientific Responses to Model--Observation Discrepancies}
\label{sec:model-observation}

Experimental feedback can change both the assessment of a candidate and the system's diagnosis of the source of uncertainty. A disagreement can concern a numerical prediction within an accepted model, the extraction of a physical quantity from a measured signal, or the assumptions defining the explanation itself. These require different subsequent investigations. BEAR's simulation-informed compression experiments, NMR strategy transfer, and the AI Theorist's optical-spectrum analysis are selected here because their reports document these distinct responses. They concern engineered structures and physical systems; their comparison does not presume that the same diagnostic operations apply to formal proofs or institutional simulations.

BEAR illustrates correction of a useful but imperfect physical prior. In crossed-barrel compression, finite-element analysis supplied predicted resilience $\widetilde U_E(x)$, and physical experiments supplied the energy $U_E(x)$ absorbed up to yield under the study's measurement convention. Rather than replace the simulation, a Gaussian process learned its discrepancy. For the different objective of toughness $U(x)$, defined as the area under the full measured force--displacement curve, the researchers instead introduced an effective-length response using predicted yield force $\widetilde F_y(x)$~\cite{asd0615}:
\begin{equation}
\delta(x)=U_E(x)-\widetilde U_E(x),\qquad
L(x)=\frac{U(x)}{\widetilde F_y(x)}.
\label{eq:bear-response}
\end{equation}
These are two study-specific response representations, not interchangeable physical objectives. Resilience and toughness were component-level energies measured in joules. The effective length expresses the displacement that would absorb energy $U$ at constant force $\widetilde F_y$; it is not a direct measurement of ductility. The elastic-response prediction did not correlate with toughness, and direct use of that prior did not accelerate simulated toughness search. Nevertheless, predicted yield force remained useful for transforming the response and downweighting designs likely to exceed the 4.8-kN test threshold~\cite{asd0615}. Thus, an approximate model can inform the scale or feasible region of a decision even when its original output poorly predicts the target. The measurement protocol helps determine which model information is useful. Researchers supplied these representations and constraints; the optimizer updated beliefs and selected designs within them. Successful discrepancy learning supports local prediction and selection without establishing a physical explanation of the residual.

NMR transfer exposes a different problem: the procedure for estimating the local physical model may cease to be appropriate. On processor Q4, two four-line multiplets were resolved, but inner transitions in the third overlapped. The original fitter required three resolved quartets. The agent adapted line assignment to infer couplings from the two resolved quartets and constrain them with the outer lines of the third; separate X- and Y-excitation spectra tested consistency~\cite{asd1064}. This preserved the three-spin model class while changing how its parameters were extracted. Importantly, a plausible fit was insufficient: none of nine initial spectral pairs passed all characterization checks, so the agent stopped and requested operator adjustment. A later acquisition passed an independent local model check before pulse optimization. This check used additional data within the same study, rather than independent-team reproduction. The scientific distinction is between fitting one spectrum and establishing a stable model suitable for subsequent control. More pulse optimization using unstable parameters would address a different uncertainty. Reuse of the adapted analysis on Q3 therefore concerned a procedure, with fresh spectra, references, and control synthesis still required.

The AI Theorist case changes explanatory assumptions rather than only a residual function or parameter-extraction procedure. For $\alpha$-RuCl$_3$, a constrained electronic starting state produced a strong optical feature absent from the observations. Releasing the orbital-occupation constraint reduced that feature and restored the observed dominant peak in the calculation. Updating quasiparticle energies while holding screening fixed changed the gap little; updating screening as well materially altered the predicted gap. Electron--hole interactions, treated through the Bethe--Salpeter equation, then produced candidate bound excitations below the calculated quasiparticle gap~\cite{asd1081}. These tests address different physical sources of disagreement: the electronic state, its screening response, and the interacting excited states. A shift of an optical peak is consequently not equivalent to a correction of the quasiparticle gap. The delegated decisions selected calculations that changed these physical assumptions, while human-acquired spectra and theoretical assessment constrained interpretation.

Agreement must also concern the observable actually measured. Reflection contrast, calculated dielectric response, and photocurrent are related but distinct quantities. Photocurrent requires charge generation and collection following excitation; matching excitation energies alone does not establish a quantitative account of that process. The study's proposed polarization contrast between candidate excitonic features supplies a further discriminating observation, but remained unmeasured~\cite{asd1081}. The model revisions therefore support a physically developed interpretation and testable predictions, while the uniqueness of the assignments and the complete photocurrent mechanism remain unresolved. Improving agreement with the observations used to guide revision cannot substitute for testing such additional consequences.

The contrast locates the delegated decision precisely: BEAR learns within researcher-designed response representations, the NMR agent revises parameter extraction with operator involvement, and the theoretical agents change explanatory calculations. A residual magnitude alone cannot select the appropriate response. Its interpretation depends on the diagnostics available to separate predictive error, measurement instability, and an inadequate physical explanation. The cases support these distinct scientific operations without demonstrating a system that autonomously controls all three.
